% =============================================================================
%  1. Definición del producto
% =============================================================================
\section{Definición del producto}

\textbf{\NombreProducto} es una plataforma SaaS (\emph{Software como Servicio}) que
integra controles de seguridad automatizados dentro de los pipelines de Integración
y Entrega Continua (CI/CD) de equipos de desarrollo de software. Sigue el enfoque
\emph{shift-left}: en lugar de revisar la seguridad al final del ciclo, la incorpora
desde la fase de escritura de código.

El producto se conecta mediante un agente o integración nativa a las herramientas que
el cliente ya usa (GitHub, GitLab, Bitbucket, Jenkins, Azure DevOps). En cada
\emph{commit} o \emph{build}, ejecuta automáticamente:

\begin{itemize}
  \item \textbf{SAST} (\emph{Static Application Security Testing}): análisis estático
        del código fuente en busca de vulnerabilidades.
  \item \textbf{SCA} (\emph{Software Composition Analysis}): detección de
        vulnerabilidades en dependencias y librerías de terceros.
  \item \textbf{Escaneo de secretos}: detecta contraseñas, llaves API y tokens
        expuestos accidentalmente en el código.
  \item \textbf{Escaneo de contenedores e IaC} (\emph{Infrastructure as Code}): revisa
        imágenes Docker y archivos de configuración (Terraform, Kubernetes).
  \item \textbf{Dashboard de cumplimiento}: tablero unificado con métricas,
        priorización de riesgos y reportes de cumplimiento (ISO 27001, PCI-DSS,
        Ley N.\textsuperscript{o} 8968 de Protección de Datos de Costa Rica).
\end{itemize}

\paragraph{Componente tecnológico.} La plataforma es 100\,\% cloud-native, con motor
de análisis basado en contenedores, una capa de inteligencia artificial para priorizar
y sugerir remediaciones de vulnerabilidades, y APIs para integración. Esto cumple con
el requisito de ser un producto con componentes tecnológicos.

\paragraph{Modelo de negocio.} Suscripción mensual/anual por número de desarrolladores
activos (modelo \emph{per-seat}), con tres planes: \textbf{Starter} (PYMEs),
\textbf{Business} (empresas medianas) y \textbf{Enterprise} (corporaciones y sector
financiero).
